\begin{lemma}[Uniform matching set-image concentration]
Fix $0<d\le1$, lower density bounds $0<\alpha,\beta\le1$, and an error $\eta>0$.  There is $\varepsilon_0>0$ such that, whenever $0<\varepsilon\le\varepsilon_0$, there are $c<1$ and $n_0$ with the following property.  If $G[U,W]$ is an $(\varepsilon,d)$-super-regular pair with $|U|=|W|=n\ge n_0$, $S\subseteq U$, $T\subseteq W$, $|S|\ge\alpha n$, $|T|\ge\beta n$, and $\sigma$ is a uniformly random perfect matching of $G$, then
\[
 \Pr\left(\left||\sigma(S)\cap T|-{|S||T|\over n}\right|>\eta n\right)\le c^n.
\]
The same $c,n_0$ work simultaneously for all permitted sizes and choices of $S,T$.
\end{lemma}
\begin{proof}
Write $s=|S|$, $t=|T|$, and $X=|\sigma(S)\cap T|$.  For a possible value $x$, a matching counted by $X=x$ splits into matchings on four balanced rectangles of orders
\[
 y_1=x,\quad y_2=s-x,\quad y_3=t-x,\quad y_4=n-s-t+x.
\]
Choosing their vertex sets first gives
\[
 M(G,x)\le {s\choose x}{t\choose x}{n-t\choose s-x}{n-s\choose t-x}
             \prod_{i=1}^4 M(y_i).                                 \tag{1}
\]
We will decrease $\varepsilon_0$ whenever necessary, in particular so that
$2\varepsilon<d$, and take $n_0$ at least the maximum of the thresholds in
\noderef{RegularInducedMatchingUpperBound} and \noderef{MatchingEnumeration}.
For $y_i\ge\sqrt\varepsilon n$, \noderef{RegularInducedMatchingUpperBound} gives
$M(y_i)\le D^{y_i}y_i!$, where $D=d+\varepsilon+3\sqrt\varepsilon$; for a smaller $y_i$ use $M(y_i)\le y_i!$.  At most three of the four $y_i$ can be small unless $n<4\sqrt\varepsilon n$, which is excluded by taking $\varepsilon_0<1/16$.  Since $D$ is bounded below by $d/2$ after decreasing $\varepsilon_0$, all omitted powers cost at most
$\min(D,1)^{-3\sqrt\varepsilon n}$.  Here are the details, including rectangles
of order zero.  Let $I=\{i:y_i<\sqrt\varepsilon n\}$ and
$q=\sum_{i\in I}y_i$.  The trivial estimate is valid also for $y_i=0$, since
the empty rectangle has exactly one matching and $0!=1$.  Moreover
$|I|\le3$ and hence $0\le q\le3\sqrt\varepsilon n$.  The product of all four
matching estimates is therefore at most
\[
 D^{n-q}\prod_{i=1}^4y_i!
 \le D^n\min(D,1)^{-3\sqrt\varepsilon n}\prod_{i=1}^4y_i! .
\]
Indeed, if $D\ge1$, then $D^{n-q}\le D^n$ and the minimum is $1$; if
$0<D\le1$, then
$D^{n-q}=D^nD^{-q}\le D^nD^{-3\sqrt\varepsilon n}$.  This also covers
$q=0$ without interpreting any missing rectangle separately.

By \noderef{MatchingEnumeration}, the denominator is at least
$(d-2\varepsilon)^n n!$.  Divide (1) by it.  The factorial cancellation can
be checked directly: substituting
$y_1=x,y_2=s-x,y_3=t-x,y_4=n-s-t+x$ into the four binomial coefficients gives
\[
 {s\choose x}{t\choose x}{n-t\choose s-x}{n-s\choose t-x}
 {\prod_i y_i!\over n!}
 ={{t\choose x}{n-t\choose s-x}\over {n\choose s}}.
\]
(If one of the $y_i$ is zero, the same identity holds because $0!=1$.)
Consequently
\[
 \Pr[X=x]\le R(\varepsilon)^n
 { {t\choose x}{n-t\choose s-x}\over {n\choose s}},
 \qquad
 R(\varepsilon):={D\over d-2\varepsilon}
                  \min(D,1)^{-3\sqrt\varepsilon}.                 \tag{2}
\]
This factor tends to $1$ as $\varepsilon\downarrow0$.  Indeed,
$D\to d$, so $D/(d-2\varepsilon)\to1$.  Also $D\ge d/2>0$, and hence
$\log\min(D,1)$ remains bounded and converges to $\log d$ (with value $0$
when $d=1$).  Thus
$-3\sqrt\varepsilon\log\min(D,1)\to0$, proving that the second factor tends
to $1$.  In particular $\log R(\varepsilon)\to0$, uniformly in $s,t,x$.

The fraction in (2) is the mass of a hypergeometric variable $Y$ obtained by
choosing a uniformly random $s$-subset of an $n$-set having $t$ marked
elements.  Put $p=t/n$ and $m=sp=st/n$.  We first prove the moment estimate
needed for both tails.  For every real $\lambda$, let $a_1,\ldots,a_n$ consist
of $t$ copies of $e^\lambda$ and $n-t$ copies of $1$.  Directly expanding over
the uniformly chosen subset gives
\[
 \mathbb E e^{\lambda Y}
 ={e_s(a_1,\ldots,a_n)\over {n\choose s}},                         \tag{3}
\]
where $e_s$ is the $s$-th elementary symmetric sum.

For completeness, we prove the elementary-symmetric mean inequality used
here.  For arbitrary nonnegative $a_1,\ldots,a_n$, replacing two entries
$a_i,a_j$ by their common average leaves every term of $e_s$ containing
neither or exactly one of them unchanged in total, while it replaces the
factor $a_i a_j$ in the remaining terms by
$((a_i+a_j)/2)^2\ge a_i a_j$.  Thus this averaging cannot decrease $e_s$ and
preserves the arithmetic mean.  Repeatedly average a largest and a smallest
entry.  The sum of squared deviations from the mean decreases, and its only
possible limit is zero: if the limiting deviation were positive, the range
would stay bounded below, so every further averaging would decrease the sum
of squared deviations by a fixed positive amount.  Continuity of $e_s$ then yields
\[
 {e_s(a_1,\ldots,a_n)\over {n\choose s}}
 \le \left({a_1+\cdots+a_n\over n}\right)^s.
\]
Applied to (3), this proves, for every real $\lambda$ (positive or negative),
\[
 \mathbb E e^{\lambda Y}\le(1-p+pe^\lambda)^s,
 \qquad
 \mathbb E e^{\lambda(Y-m)}
 \le \exp\bigl(m(e^\lambda-1-\lambda)\bigr),                      \tag{4}
\]
where the second inequality uses $\log(1+z)\le z$.

We now derive the two tails separately.  For $\lambda>0$, Markov's inequality
and (4) give
$\Pr(Y\ge m+u)\le
\exp(m(e^\lambda-1-\lambda)-\lambda u)$.  If $m>0$, choosing
$\lambda=\log(1+u/m)$ gives exponent $-m h(u/m)$, where
$h(z)=(1+z)\log(1+z)-z$.  The elementary inequality
$h(z)\ge z^2/(2(1+z/3))$ for $z\ge0$ follows by differentiating the
difference twice: its second derivative is
$1/(1+z)-27/(z+3)^3\ge0$, since
$(z+3)^3-27(z+1)=z^3+9z^2\ge0$, and both its value and first derivative
vanish at zero.  This gives the first bound below.  If $m=0$, then $Y=0$ and the assertion
is immediate.  For the lower tail take $\lambda<0$.  Since
$e^\lambda-1-\lambda\le\lambda^2/2$ for $\lambda\le0$ (compare second
derivatives, using $e^\lambda\le1$, and the common value and derivative at
zero), Markov applied to
$e^{\lambda(Y-m)}$ and the choice $\lambda=-u/m$ give the second bound when
$m>0$; again $m=0$ is immediate.  We have proved, for $u\ge0$,
\[
 \Pr(Y\ge m+u)\le e^{-u^2/(2(m+u/3))},\qquad
 \Pr(Y\le m-u)\le e^{-u^2/(2m)}.                                 \tag{5}
\]

Sum (2) over the two tails $|x-st/n|>\eta n$.  The hypergeometric masses in
(2) sum to the corresponding probabilities, so no additional count of
possible values is needed.  Since $m\le n$, (5) gives
\[
 \Pr(|X-st/n|>\eta n)
 \le 2\exp\{n\log R(\varepsilon)-\eta^2n/(2(1+\eta/3))\}.         \tag{6}
\]
Choose $\varepsilon_0$ so that
$\log R(\varepsilon)<\eta^2/(4(1+\eta/3))$ throughout
$0<\varepsilon\le\varepsilon_0$.  Then (6) is at most
$2e^{-\gamma n}$ with $\gamma=\eta^2/(4(1+\eta/3))>0$.
Enlarge $n_0$ so that $2e^{-\gamma n}\le e^{-\gamma n/2}$ for
$n\ge n_0$, and take $c=e^{-\gamma/2}<1$.  Every choice made depends only on
$d$ and $\eta$ (and hence is certainly allowed to depend on
$d,\alpha,\beta,\eta$), not on $S,T,s,t$; this proves the asserted
simultaneous uniformity.  The lower bounds involving $\alpha$ and $\beta$ are
therefore harmless restrictions on the sets, rather than ingredients in the
large-rectangle threshold.
\end{proof}
